\documentclass[10pt,a4paper]{amsart}
\usepackage[margin=19mm]{geometry}
\usepackage{amsmath,amssymb,mathtools}
\usepackage[scr=rsfs,cal=euler]{mathalfa}
\usepackage{xfrac}
\usepackage[unicode,hidelinks,bookmarks=false]{hyperref}
\newcommand{\ie}{i.e.\ }
\newcommand{\cf}{cf.\ }
\newcommand{\resp}{respectively\ }
\newcommand{\Subsection}{\subsection}
\providecommand{\nobreakdash}{\nobreak\hbox{-}}
\setlength{\emergencystretch}{2em}
\multlinegap0pt
\usepackage{bm}
\usepackage{bbm}
\usepackage{dsfont}
\usepackage{xspace}
\usepackage{cancel}
\usepackage{mathtools}
\usepackage{amsthm}
\makeatletter
\@ifundefined{theorem}{\newtheorem{theorem}{Theorem}}{}
\@ifundefined{lemma}{\newtheorem{lemma}[theorem]{Lemma}}{}
\makeatother
\DeclareMathOperator{\F}{gF}
\DeclareMathOperator{\gid}{\textnormal{gId}}
\DeclareMathOperator{\var}{var}
\newcommand{\cV}{\mathcal{V}}
\title[Cocharacters of generalized polynomial identities]{Cocharacters of generalized polynomial identities}
\author{Sebastiano Argenti, Giovanni Busalacchi}
\date{Source: arXiv:2508.00464v1 (CC BY 4.0)}
\begin{document}
\maketitle

\noindent\emph{Source and licence.} Cocharacters of generalized polynomial identities. Version arXiv:2508.00464v1 (CC BY 4.0), distributed under the Creative Commons Attribution 4.0 International licence, as recorded in the arXiv licence metadata for this exact version and shown on the abstract page https://arxiv.org/abs/2508.00464v1. Full rights notice: licenses/arxiv-2508.00464-CC-BY-4.0.txt.\par\smallskip

\noindent\textit{Editorial scope.} Counted unit: the Theorem whose source label is \texttt{th: Grassmann envelope} in the article (source0.tex lines 593--596, source label \texttt{th: Grassmann envelope}), delivered together with the article's own complete original proof of it (source0.tex lines 597--627, one \texttt{proof} environment of 31 lines). No mathematics is written, repaired or re-derived: statement and proof are the article's own text line for line. Required context retained uncounted: 3.7.4 (lines 554--560); 2.5.6 (lines 565--580); 2.4 (lines 568--571). Every definition, display and label of those lines is retained unchanged. Omitted from this package and not redistributed: 1--553, 561--564, 572--592, 628--710. Nothing inside a retained excerpt is omitted or altered: every retained body is delivered exactly as the article's lines are written, so no omission receipt of any kind accompanies this package. Citations: the retained excerpts carry no citation command at all, so no bibliography entry of the article is transcribed and no reference is redistributed. Reference adaptation: none was needed and none was made; every reference inside the delivered unit resolves inside it, so no unresolved reference or citation is produced. Printed slips kept literal: no printed word, symbol, hypothesis or number of the counted statement or of its proof was altered. Layout and class: the article's own class is \texttt{article} with packages including graphicx, amsmath, amsfonts, amssymb, amsthm, marvosym, xcolor, comment, fontenc, inputenc, babel, tikz-cd, fullpage; the wrapper is the lane's pinned \texttt{amsart} editorial wrapper at 10pt on A4, which restates the theorem-like environments the retained text needs and adds the article's own macro definitions that the retained text needs (\texttt{\textbackslash F}, \texttt{\textbackslash cV}, \texttt{\textbackslash gid}, \texttt{\textbackslash var}), with extra wrapper packages (bm, bbm, dsfont, xspace, cancel, mathtools). The delivered numbering is the unit's own. Assets: no retained span contains a figure environment, a graphics-inclusion command, a PDF-page inclusion, a table or a TikZ picture; no image file is copied into this package, the package declares no asset dependency and no image file is redistributed with it. Licence: version arXiv:2508.00464v1 of this article is distributed under the Creative Commons Attribution 4.0 International licence, as recorded in the arXiv licence metadata for this exact version and shown on the abstract page https://arxiv.org/abs/2508.00464v1. A full rights notice is delivered at licenses/arxiv-2508.00464-CC-BY-4.0.txt. Characters: every retained excerpt is pure ASCII, so the delivered engine, XeLaTeX with its default Latin Modern text font, reports no missing character.
\begin{lemma}\label{3.7.4}
Let $f\in gP_{l,m}$. Then: 
  \begin{enumerate}
    \item [1)] $\tilde{\tilde{f}}=f$;
    \item [2)] $f$ is a graded generalized identity of $E(A)$ if and only if $\Tilde{f}$ is a graded generalized identity of $A$.
  \end{enumerate}
\end{lemma}

\begin{lemma}
    \label{2.5.6}
    Let $\lambda \vdash n, T_\lambda$ a Young tableau of shape $\lambda$ and $f=e_{T_\lambda}g$ where $g=g(x_1,\ldots,x_n)$ is some multilinear generalized polynomial in $x_1,\ldots,x_n$. If $\lambda \in H(k,l)$, then there exists a decomposition of $X^{(n)}=\{ x_1,\ldots,x_n\}$ into a disjoint union
    \begin{equation}
        \label{2.4}
        X^{(n)}=X_1 \cup \ldots \cup X_{k'}\cup Y_1 \cup \ldots \cup Y_{l'}
    \end{equation}
    with $k' \le k, l' \le l$ and a multilinear polynomial $f'=f'(x_1,\ldots,x_n)$ such that 
    \begin{itemize} 
        \item $f'$ is symmetric on the variables in each set $X_i, 1 \le i \le k'$;
        \item $f'$ is alternating on the variables in each set $Y_j, 1 \le j \le l'$;
        \item $FS_nf=FS_nf'$;
        \item the integers $k',l', |X_1|,\ldots,|X_{k'}|,|Y_1|,\ldots,|Y_{l'}|$ are uniquely determined by $\lambda$ and do not depend on the choice of the tableau $T_\lambda$;
        \item the decomposition \eqref{2.4} is uniquely defined by $T_\lambda$ and does not depend on $g$.
    \end{itemize}
\end{lemma}

\begin{theorem}
\label{th: Grassmann envelope}
    Let $W$ be a finite dimensional associative unital algebra and let $A$ be a $W$-algebra whose generalized cocharacter multiplicities are contained in the hook $H(k,l)$. Then, there exists a finitely generated superalgebra $L=L^{(0)} \oplus L^{(1)}$ with $k$ even generators and $l$ odd generators such that $\gid(A)=\gid(E(L))$.
\end{theorem}

\begin{proof}
    Let $\cV=\var^W(A)$ and consider the variety of $W$-superalgebras $\mathcal{V}^*$ and its relatively free superalgebra 
    \[ L = \F_{k,l}(\cV^*)=\frac{W\langle y_1,\dots,y_k,z_1,\dots,z_l\rangle}{W\langle y_1,\dots,y_k,z_1,\dots,z_l\rangle\cap \gid_2(\cV^*)} \]
    on $k$ even free generators $l$ odd free generators. We shall prove that $\mathcal{V}$ is generated by $E(L)$.

    Clearly, $E(L) \in \mathcal{V}$ by the definition of $\mathcal{V}^*$. We next verify that $A$ satisfies all identities of $E(L)$.

    Let $f \in \gid(E(L))$ be a multilinear identity of $G(L)$ of degree $n$. We can assume that $f=e_{T_\lambda}g$ for some multilinear polynomial $g\in gP_n$. If $\lambda \notin H(k,l)$, then $f$ is an identity of $A$ since $g\chi_n(A)$ lies in the hook $H(k,l)$. Now, let $\lambda \in H(k,l)$. By applying Lemma \ref{2.5.6}, we may assume that
    \[
    f=f(y_1^1,\ldots,y_1^{m_1},\ldots,y_k^1,\ldots,y_k^{m_k},z_1^1,\ldots,z_1^{r_1},\ldots,z_l^1,\ldots,z_l^{r_l})
    \]
    depends on $k$ (possibly empty) sets of symmetric variables $\{ y_i^1,\ldots,y_i^{m_i} \}, 1 \le i \le k$ and on $l$ (possibly empty) sets of alternating variables $\{z_j^1,\ldots,z_j^{r_j}\}, 1 \le j \le l$. If we consider all $y_i^j$ as even variables and all $z_i^j$ as odd variables, then $f$ can be considered as a graded identity of the superalgebra $G(L)$. So, by Lemma \ref{3.7.4}, the superalgebra $L$ satisfies the graded identity $\Tilde{f}$. Now, we substitute all the variables in the set $\{ y_i^1,\ldots,y_i^{m_i} \}$ with the even generator $y_i\in L$, $1 \le i \le k$, and all the variables in the set $\{z_j^1,\ldots,z_j^{r_j}\}$ with the odd generator $z_j\in L$, $1 \le j \le l$. We obtain that \begin{equation} \label{eq: f=0 in L}
        \Tilde{f}(y_1,\ldots,y_1,\ldots,y_k,\ldots,y_k,z_1,\ldots,z_1,\ldots z_l,\ldots, z_l) = 0
    \end{equation}
    in $L$.

    Consider the superalgebra $A\otimes E = (A\otimes E_0) \oplus (A\otimes E_1)$ where the grading is determined by $E$. We note that $E(A\otimes E)= A\otimes R$ where $R=(E_0\otimes E_0) \oplus (E_1 \otimes E_1)$ is a commutative algebra, so $E(A\otimes E)$ satisfies all multilinear identities of $A$. Therefore, $E(A\otimes E)\in \cV$ and, by the definition of $\cV^*$, $A\otimes E\in \cV^*$. Now take some arbitrary elements $\Bar{y}_1^1,\ldots,\Bar{y}_1^{m_1}, \ldots,\Bar{y}_k^1,\ldots,\Bar{y}_k^{m_k}, \Bar{z}_1^1,\ldots,\Bar{z}_1^{r_1}, \ldots, \Bar{z}_l^1,\ldots,\Bar{z}_l^{r_l}$ in $A$, and consider 
    \begin{align*}
        q_i=\Bar{y}_i^1 \otimes a_i^1+\cdots+\Bar{y}_i^{m_i}\otimes a_i^{m_i}, \quad i=1,\ldots,k\\
        p_j=\Bar{z}_j^1\otimes b_j^1+\cdots+\Bar{z}_j^{r_j} \otimes b_j^{r_j}, \quad j=1,\ldots,l
    \end{align*}
    where $a_1^1,\ldots,a_k^{m_k} \in E_0$, $b_1^1,\ldots,b_l^{r_l}\in E_1$ are monomials of the Grassmann algebra $E$ written on distinct generators of $E$, hence $a_1^1\cdots a_k^{m_k}b_1^1\cdots b_l^{r_l} \neq 0$.

    Since $L$ is a relatively free algebra of $\cV^*$ and $A\otimes E\in \cV^*$, there exists an homomorphism of $W$-superalgebras sending $y_i\in L$ in $q_i\in A\otimes E$, $1 \le i \le k$, and $z_j\in L$ in $p_j$, $1 \le j \le l$. Applying such homomorphism to Equation \ref{eq: f=0 in L}, we get
    \begin{align*}
    0=\Tilde{f}(q_1,\ldots,p_l)=&\\
    m_1!\cdots m_k!r_1!&\cdots r_l!\ f(\Bar{y}_1^1,\ldots,\Bar{y}_1^{m_1},\ldots,\Bar{z}_l^1,\ldots,\Bar{z}_l^{r_l}) \otimes a_1^1\cdots a_1^{m_1}\cdots b_l^1 \cdots b_l^{r_l},
 \end{align*}
 since $(a_i^j)^2=(b_i^j)^2=0$ for all $i,j$, and $\Tilde{f}$ is symmetric on any set $\{ y_i^1,\ldots,y_i^{m_i} \}$ and on any set $\{ z_j^1,\ldots,z_j^{r_j} \}$.
 Therefore, $f$ vanishes on every evaluation of the variables with elements of $A$, that is $f\in \gid(A)$. We conclude that $\gid(E(L))\subseteq \gid(A)$ and we are done.
\end{proof}

\end{document}
